
\section{Quelques généralités}


\begin{frame}
\frametitle{Systèmes distribués}

Un \alert{système distribué} est une application qui coordonne les
actions de plusieurs ordinateurs connectés par un réseau.

  
  \vfill
  La coordination consiste à échanger des \alert{messages}, codés selon un certain
    \alert{protocole}.
      
   $\Rightarrow$ on se limite ici aux architectures ``\textit{shared-nothing}''  -> 
       pas de mémoire partagée, pas de disque partagé, etc.
 
   
   \vfill 
   
   Le réseau se charge de la transmission des messages, selon un mode qui dépend
      de l'infrastructure sous-jacente.
    
   $\Rightarrow$ Local area networks (LAN), Peer to peer (P2P) networks, \ldots
    
    \vfill
    
    Le système est constitué de \alert{n{\oe}uds}\,: des composants logiciels
      capables d'émettre/recevoir des messages selon le protocole.
      
\begin{remark}
C'est la présence d'un \alert{composant logiciel} qui définit le rôle d'une machine.
Une même machine peut être simultanément (multi-)serveur(s) et (multi-)client(s)
\end{remark}
\end{frame}

\begin{frame}
\frametitle{Un exemple de système distribué\,: le Web}

Il obéit à un protocole\,: lequel\,?

\vfill

Il s'appuie sur un réseau\,: lequel?

\vfill

Il donne accès à des resources, grâce à des \alert{clients web}. Par exemple\,:?

\vfill

Les ressources sont gérées par des \alert{serveurs}. Par exemple\,: ?

\vfill

Basé sur un système d'adressage des ressources\,: ?

\vfill

Connaissez vous une machine qui  est multi-serveurs et multi-clients
(regardez bien autour de vous!)\,,
\end{frame}



\begin{frame}
\frametitle{Applications analytiques (par lot) distribuées}

Exemple d'une \alert{application analytique} (lecture de masses de données).

\smallskip

% A two-columns presentation
\begin{tabular}{l|p{6cm}}
   \parbox{5cm}{

  \textbf{Lecture séquentielles.}  Il faut 166 minutes (plus de deux heures et demi) pour lire un Téra-octet 
  
  \vfill
  \textbf{Lecture parallèle} Avec 100 disques, en supposant une lecture séquentielle et en parallèle: environ 1mn 30s.
 \vfill
 
 \textbf{Lecture distribuée} Avec 100 machines, chacune dotée d'un disque, d'un CPU et d'un fragment des données (partitionnement),
     1 mn 30s, sans goulot d'étrangement pour le traitement. 
   }  % Inclusion of an XML file
  & 
  \parbox{6cm}{\figSlideWithSize{bigpic-distrstorage}{5.5cm} }
\end{tabular}
\vfill

% A Beamer box
\begin{beamerboxesrounded}{Scalabilité horizontale}
La distribution augmente le débit en proportion du nombre de disques.
\end{beamerboxesrounded}

\end{frame}

\begin{frame}
\frametitle{Applications transactionnelles (temps réel) distribuées}

Exemple d'une \alert{application transactionnelle} (100 lectures aléatoires).

\smallskip


% A two-columns presentation
\begin{tabular}{l|p{6cm}}
   \parbox{5cm}{

  \textbf{Une seule machine, données sur disque}. 100 lectures, 10 ms par lecture, env. 1 s.
  
  \vfill
  \textbf{Partionnement sur disque}. Dans le meilleur des cas, chaque lecture sur un disque\,:
     10 ms.
 \vfill
 
 \textbf{Partitionnement en mémoire}. 100 lectures en RAM: quelques micro-secondes. 
   }  % Inclusion of an XML file
  & 
  \parbox{6cm}{\figSlideWithSize{bigpic-transaction}{5.5cm} }
\end{tabular}
\vfill

% A Beamer box
\begin{beamerboxesrounded}{Scalabilité horizontale}
La distribution vise à placer les données en mémoire, au moins pour les lectures.
\end{beamerboxesrounded}

\end{frame}


\begin{frame}
\frametitle{Caractéristiques des systèmes distribués}

Sujet complexe... On va se concentrer sur les aspects essentiels suivants.

\begin{itemize}
  \item la \alert{réplication des données}, indispensable pour assurer
     
  \begin{itemize}
    \item  la \alert{tolérance aux pannes}. Plus il y a de serveurs, plus le risque de panne est élevé\,;
         il faut donc maintenir les données disponibles dans tous les cas.
     \item les \alert{performances}. La réplication peut \alert{aussi} assurer 
      (partiellemnt) le passage à l'échelle.
  \end{itemize}
 
 \item le \alert{partionnement} (\textit{sharding}), autre outil, plus puissant, du passage à l'échelle
  \begin{itemize}
    \item  vous avez dix fois plus de documents\,? mettez 10 fois plus de machines, et 1/10ème des
        documents sur chaque machine.
     \item comment trouver la machine qui héberge la donnée qui vous intéresse\,?
  \end{itemize}

\item les \alert{traitements par lot}, typique des opérations analytiques (OLAP).
  \begin{itemize}
    \item  la distribution doit bénéficier aux \alert{accès aux données} (cf. ci-dessus)
     \item les \alert{traitements} doivent eux-aussi être distribués.
  \end{itemize}
   
\end{itemize}
\end{frame}


\begin{frame}
\frametitle{Caractérisation des applications (et des systèmes qui vont bien)}

\alert{Applications transactionnelles}. Exemple\,: e-commerce.

 \begin{itemize}
    \item  contraintes temps réel\,: les requêtes doivent être très rapides (quelques ms).
    \item une transaction accède à une proportion très faible des données.
    \item lectures/écritures équilibrées (dépend de l'appli.)
    \item besoins d'une \alert{cohérence} plus ou moins forte.
     \item \alert{Systèmes NoSQL}\,: \mongodb/, \couchdb/, Voldemort, Amazon S3, et beaucoup d'autres.
  \end{itemize}

\alert{Applications analytiques} Exemple\,: classification, \textit{clustering}, stats.

 \begin{itemize}
    \item pas de contrainte temps réel (certains traitements durent des semaines!)
    \item une transaction peut accèder à \alert{toutes} les données.
    \item beaucoup de lectures, très peu d'écritures.
    \item pas de notion de  \alert{cohérence} (transactionnelle).
     \item \alert{Systèmes}\,: \textsc{Hadoop}, \mapreduce/
  \end{itemize}
\end{frame}

\end{document}


